\section{Practical rules, and what this term's paper does \emph{not} claim}
\label{sec:scope}

\emph{Source: handout \S1.7--1.9, \S2. Status: skeleton from a first read.}

\subsection{Two cheap but important rules (\S1.7)}

\newtag, cheap, important.
\begin{enumerate}
  \item \textbf{Explanatory data must be held out.} $Y_i^* = f(X_i)$ on rows
    used to \emph{train} $f$ is not a fair sample of what $f$ does (a random
    forest nearly reproduces its own training outcomes). Draw every curve
    on rows the model never saw. Out-of-bag / cross-fitted predictions are
    predictions of \emph{different} models, not of $f$ itself --- don't
    substitute them.
  \item \textbf{Name the reference population.} The outer average is over
    some distribution of $W$; training cohort, held-out sample, and
    deployment population give genuinely different curves whenever the
    effect varies with $w$. This term's inference is all for the
    \emph{held-out sample's} population. Carrying a curve to a different
    population (``transport'') needs extra data and different formulas ---
    explicitly \textbf{out of scope} this term.
\end{enumerate}

\subsection{What Paper A does \emph{not} claim (\S1.8)}

Not claiming: that parent adjustment, AIPW, cross-fitting, or dose-response
bands are new ideas; that the method works ``without a causal model'' (it
always needs the ECM); that this is the first semiparametric inference for
causal explainability (cf.\ Zhang \& Gao 2026; Gao \& Zhao 2024 --- different
quantities, worth knowing about but not the same target); bands for
\emph{individual per-person} curves (out of scope --- only aggregate curves
here).

\textbf{Why this matters for us as interns:} knowing the boundary of the
claim keeps write-ups honest and stops accidental overclaiming of novelty
that belongs to cited work --- exactly what the REUSE/NEW tagging throughout
the handout is for.

\subsection{Paper C --- later project, not this term (\S2)}

A theory paper making the \S\ref{sec:bigpicture} picture rigorous: why a
\emph{deterministic} output node removes the cross-world content of edge
interventions; influence functions for all four plots (NIDP's is Fulcher's;
PCDP's is a g-formula one); efficiency results when the whole DAG is used;
and \textbf{graph uncertainty} --- from observational data one often cannot
orient every edge, so there are several candidate parent sets and several
candidate curves. Caution flagged in the handout: for linear-Gaussian
models, all graphs in a Markov-equivalence class have identical likelihood,
so a BIC ``neighbourhood'' of orientations is \emph{empty} there; for
nonlinear models, orientations actually get scored differently. A couple of
intern-sized pieces exist here (simulating candidate curves with IDA
software; a merging algorithm) but the core theory is research-level ---
\textbf{not this term's scope}, noted here mainly so this term's work
doesn't accidentally drift into it.

\begin{openquestions}
  \item None yet --- revisit once Paper A's early tasks
    (\S\ref{sec:identification}--\S\ref{sec:estimation}) are underway and it
    becomes clearer where the term's actual scope will land.
\end{openquestions}
